\begin{helper}
Let $G$ be a finite graph, let $\mu$ be an activation--priority sample with
activation probability $\gamma/\Delta$ as in
\noderef{RandomIndependentSetSampling}, and assume $\Delta>0$.  Fix three
distinct vertices $a,b,c$, put $Q=\{a,b,c\}$, and consider the ordered chamber
\[
0\le z\le y\le x\le \gamma
\]
corresponding to $x=\gamma(1-\pi(a))$, $y=\gamma(1-\pi(b))$, and
$z=\gamma(1-\pi(c))$.  For an activation set $A\supseteq Q$, let the three
atomwise exponents be
\[
|\{w\in A\setminus Q:w\sim a\}|,
\quad
|\{w\in A\setminus Q:w\not\sim a,\ w\sim b\}|,
\quad
|\{w\in A\setminus Q:w\not\sim a,\ w\not\sim b,\ w\sim c\}|.
\]
Then the finite sum over activation atoms of the Bernoulli atom weight times
the corresponding ordered-chamber integral is exactly
\[
\begin{aligned}
&\sum_{A\supseteq Q}
\left(\frac{\gamma}{\Delta}\right)^{|A|}
\left(1-\frac{\gamma}{\Delta}\right)^{|V(G)|-|A|}
\frac1{\gamma^3}
\int_{0\le z\le y\le x\le\gamma}
\left(1-\frac{x}{\gamma}\right)^{|\{w\in A\setminus Q:w\sim a\}|} \\
&\qquad\qquad\qquad\qquad\cdot
\left(1-\frac{y}{\gamma}\right)^{|\{w\in A\setminus Q:w\not\sim a,\ w\sim b\}|}
\left(1-\frac{z}{\gamma}\right)^{|\{w\in A\setminus Q:w\not\sim a,\ w\not\sim b,\ w\sim c\}|}
\,dx\,dy\,dz  \\
&=
\frac1{\Delta^3}
\int_{0\le z\le y\le x\le\gamma}
\left(1-\frac{x}{\Delta}\right)^{|\{w\in V(G)\setminus Q:w\sim a\}|}
\left(1-\frac{y}{\Delta}\right)^{|\{w\in V(G)\setminus Q:w\not\sim a,\ w\sim b\}|} \\
&\qquad\qquad\qquad\qquad\cdot
\left(1-\frac{z}{\Delta}\right)^{|\{w\in V(G)\setminus Q:w\not\sim a,\ w\not\sim b,\ w\sim c\}|}
\,dx\,dy\,dz .
\end{aligned}
\]
\end{helper}
\begin{proof}
The activation law in \noderef{RandomIndependentSetSampling} says that an
activation atom $A$ has weight
\[
\left(\frac{\gamma}{\Delta}\right)^{|A|}
\left(1-\frac{\gamma}{\Delta}\right)^{|V(G)|-|A|}.
\]
Since $\Delta>0$, \noderef{RandomIndependentSetSamplingGammaLeDelta} gives
$\gamma\le\Delta$ in every nonempty sampling instance, and the definition of
the sampling law gives $\gamma>0$.  Thus $\gamma/\Delta$ is a genuine
Bernoulli parameter and the change of variables
$x_q=\gamma(1-\pi(q))$ contributes the density factor $\gamma^{-3}$ for the
three fixed coordinates $a,b,c$.

Fix chamber coordinates $(x,y,z)$ with $0\le z\le y\le x\le\gamma$.  Because
the sum over $A$ is finite and the chamber integrands are polynomials in
$x,y,z$, the finite sum may be moved inside the three integrals.  We therefore
evaluate the activation sum pointwise at this fixed triple of coordinates.

The condition $A\supseteq Q$ forces $a,b,c$ to be activated.  Their three
Bernoulli activation factors are
\[
\left(\frac{\gamma}{\Delta}\right)^3.
\]
After multiplication by the priority-density factor $\gamma^{-3}$, these
three forced activations give the outside coefficient $\Delta^{-3}$.

Now consider any vertex $w\notin Q$.  Its contribution to the pointwise
activation sum depends only on the first applicable class in the ordered list
\[
w\sim a,\qquad
w\not\sim a\text{ and }w\sim b,\qquad
w\not\sim a,\ w\not\sim b\text{ and }w\sim c,
\]
or on none of these classes.  If $w\sim a$, then an inactive $w$ contributes
$1-\gamma/\Delta$, while an active $w$ contributes
$(\gamma/\Delta)(1-x/\gamma)$.  The total contribution of the two activation
choices is therefore
\[
1-\frac{\gamma}{\Delta}
+\frac{\gamma}{\Delta}\left(1-\frac{x}{\gamma}\right)
=1-\frac{x}{\Delta}.
\]
If $w\not\sim a$ and $w\sim b$, the same calculation gives
$1-y/\Delta$.  If $w\not\sim a$, $w\not\sim b$, and $w\sim c$, it gives
$1-z/\Delta$.  If $w$ is in none of the three classes, no chamber survival
factor involves $w$, and its inactive and active contributions add to
\[
1-\frac{\gamma}{\Delta}+\frac{\gamma}{\Delta}=1.
\]

The activation choices of distinct vertices are independent and the
activation-set sum ranges over all choices with $Q$ forced active.  Hence the
pointwise sum factors as the product of the preceding one-vertex
contributions over all vertices outside $Q$.  Collecting like factors gives
\[
\left(1-\frac{x}{\Delta}\right)^{|\{w\in V(G)\setminus Q:w\sim a\}|}
\left(1-\frac{y}{\Delta}\right)^{|\{w\in V(G)\setminus Q:w\not\sim a,\ w\sim b\}|}
\left(1-\frac{z}{\Delta}\right)^{|\{w\in V(G)\setminus Q:w\not\sim a,\ w\not\sim b,\ w\sim c\}|}.
\]
Together with the forced-activation coefficient $\Delta^{-3}$, this is the
integrand on the right-hand side.  Integrating the pointwise identity over the
ordered chamber proves the displayed equality.
\end{proof}
